% Cuestión IX. Dificultades de implementación. Cada punto remite a un artefacto del repositorio
% (docs/LSBI-NOTES.md, docs/CRYPTO-NOTES.md, un caso de prueba, el Makefile, .gitattributes o los resúmenes de fase).
% Los valores citados se volvieron a ejecutar antes de escribirlos.

Las dificultades que siguen son las que quedaron registradas en el repositorio durante la implementación, no una reconstrucción de memoria. Cada una nombra el archivo donde puede
comprobarse: las notas \texttt{docs/LSBI-NOTES.md} (entradas A1 a A10) y \texttt{docs/CRYPTO-NOTES.md} (entradas K1 a K8), un caso de \texttt{tests/run\_tests.sh}, el \texttt{Makefile}, el
archivo \texttt{.gitattributes} o los resúmenes de cada fase. Las ambigüedades del propio documento de Majeed y Sulaiman~\cite{majeed2015} se analizan en la \hyperref[sec:cuestion-i]{Cuestión I}; aquí se cuenta
cómo se enfrentaron al programar.

\subsection{El algoritmo LSBI}

El paper deja sin definir casi todo lo que hace falta para escribir código: en qué orden se recorren los bytes, qué canales llevan datos, dónde se guarda el registro de patrones
invertidos, sobre qué bytes se cuentan los cambios, qué pasa en un empate y en qué orden se leen los dos bits del patrón. Con varias lecturas posibles para cada una de esas preguntas, escribir
directamente el programa en C habría impedido distinguir un error de lectura de un error de código. Por eso primero se escribió un modelo de referencia en Python
(\texttt{tools/lsbi.py}) con un interruptor por cada lectura alternativa, y se lo enfrentó a los tres archivos LSBI de la cátedra antes de escribir una línea de C.

La sonda (\texttt{python3 tools/lsbi.py probe}) decodifica el vector con la lectura elegida e informa el tamaño leído, si el encuadre es válido y si volver a insertar el mensaje reproduce el archivo. El cuadro~\ref{tab:q9-hipotesis}
resume lo que dio cada lectura sobre \texttt{Ejemplo/lado.bmp} y el vector correspondiente. El tamaño correcto es 44\,886 bytes (la imagen PNG oculta). Seis lecturas se descartan y dos no pueden decidirse con los vectores.

\begin{table}[htbp]
\centering
\footnotesize
\renewcommand{\arraystretch}{1.05}
\begin{tabularx}{\textwidth}{@{}>{\raggedright\arraybackslash}p{3.2cm}>{\raggedright\arraybackslash}p{3.6cm}>{\raggedright\arraybackslash}X>{\raggedright\arraybackslash}p{1.9cm}@{}}
\toprule
Lectura alternativa & Interruptor de la sonda & Resultado observado & Conclusión \\
\midrule
Los bytes rojos también llevan mensaje & \texttt{--use-red} & tamaño 658, encuadre de texto cifrado, la reinserción no coincide & Descartada \\
El salto del rojo se cuenta desde el byte 4 & \texttt{--red-relative} & tamaño 109\,078, encuadre inválido, la reinserción no coincide & Descartada \\
Bandera del patrón $p$ en el byte $3-p$ & \texttt{--flag-order reversed} & sobre el vector AES-256 (banderas \texttt{0001}): tamaño 4\,294\,922\,424, encuadre inválido, la reinserción no coincide & Descartada \\
Bits de cada byte del mensaje del menos al más significativo & \texttt{--bit-order lsb} & tamaño 62\,826, encuadre de texto cifrado (pero la reinserción \emph{sí} coincide) & Descartada \\
Los cambios se cuentan también en los bytes rojos & \texttt{--count-red} & banderas recalculadas \texttt{0000} contra \texttt{1111}; la reinserción no coincide & Descartada \\
Una bandera en 0 significa patrón invertido & \texttt{--flag-zero-inverts} & tamaño 4\,294\,922\,409, encuadre inválido, la reinserción no coincide & Descartada \\
Un empate de cambios invierte el patrón & \texttt{--tie invert} & la reinserción coincide igual & Indecidible \\
Los dos dígitos del patrón se leen al revés (\texttt{01} y \texttt{10}) & \texttt{--pattern-bits 12} & la reinserción coincide igual en los tres vectores & Indecidible \\
\bottomrule
\end{tabularx}
\caption{Lecturas alternativas de LSBI probadas contra los vectores de la cátedra (entradas A1 a A6 de \texttt{docs/LSBI-NOTES.md})}
\label{tab:q9-hipotesis}
\end{table}

Una lección de este trabajo está en la fila del orden de los bits: \texttt{reembed\_identical=yes} \emph{no} prueba que la decodificación sea correcta. Una decodificación equivocada también se reinserta sin
pérdida, porque codificar y decodificar con la misma regla equivocada es una identidad. Lo que discrimina es el tamaño decodificado y el hash del PNG extraído
(nota A3 y resumen de la fase 2). Por eso las notas registran en cada caso el tamaño decodificado y el \texttt{sha256} de la imagen extraída.

Lo que queda sin decidir es importante para los archivos de estegoanálisis: la regla de empate (nota A6), el orden de los dígitos del patrón (A4, que sólo importaría si las banderas de \texttt{01} y \texttt{10}
difirieran, cosa que ninguno de los tres vectores muestra), la capacidad máxima y el texto del mensaje de error (A8) y el conteo del salto del rojo cuando hay relleno de filas (A9, porque \texttt{lado.bmp}
mide 640 píxeles de ancho y no tiene relleno). Las cuatro decisiones están marcadas \texttt{flagged-unverified} en las notas, para revisarlas primero si aparece un archivo de la cátedra que las contradiga.

\subsection{Encripción}

\textbf{Una errata en el archivo de ejemplo.} \texttt{Ejemplo/README.txt} lista, para el archivo \texttt{ladoLSB1aes128cbc\allowbreak.bmp}, el vector de inicialización \texttt{212420edc583a686a94d19a3497363a2}, y una
clave derivada rotulada ``32 bytes'' que en realidad tiene 32 dígitos hexadecimales, es decir, 16 bytes. Con la derivación descrita en la consigna, \texttt{python3 tools/cryptoref.py kdf aes128 margarita}
da la misma clave (\texttt{03db0a157acfe8de523760aa731d8122}) pero el vector \texttt{b25f8d99f3173ec0b52849f459a4c20d}; el de la lista es el de AES-256. Con el vector de la lista, la lectura
del archivo falla con \texttt{decrypted size field says 2474312226 bytes but the plaintext has only 44895 bytes}, porque el primer bloque CBC sale corrupto (sólo cambia el IV, así que el relleno del último
bloque sigue siendo válido). Con el vector derivado se recupera el PNG de 44\,886 bytes. El archivo mismo decidió cuál era el error (nota K3, casos \texttt{cryptoref-kdf-aes128} y
\texttt{cryptoref-readme-aes128-iv-erratum}).

\textbf{Una sola derivación, primero la clave.} La consigna no dice cómo se separan clave y vector según el algoritmo. Los tres vectores de la cátedra obedecen una sola llamada a PBKDF2~\cite{rfc8018}
que produce $\mathit{long\_clave}+\mathit{long\_iv}$ bytes, con la clave primero (nota K2). Esa regla es la que fija AES-192, que no tiene archivo de ejemplo (nota K4, sin verificar contra un vector de la cátedra: el
oráculo independiente aplica la misma regla, de modo que prueba consistencia y no corrección).

\textbf{CFB de 8 bits y OFB de 128 bits.} Los nombres de la consigna deben traducirse con cuidado a los de OpenSSL~\cite{opensslevp}: CFB es \texttt{EVP\_*\_cfb8} en todos los algoritmos y OFB es el \texttt{EVP\_*\_ofb} de bloque
completo, con lo que el OFB de 3DES tiene forzosamente realimentación de 64 bits~\cite{nist80038a}. El vector de 3DES sólo se descifra con \texttt{cfb8} (nota K6).

\textbf{Relleno y capacidad.} Sólo ECB y CBC rellenan (PKCS5). El programa necesita conocer la longitud del texto cifrado \emph{antes} de comprobar la capacidad del portador, así que la calcula: es la del texto plano redondeada
hacia arriba al siguiente múltiplo del bloque (más un bloque si ya era múltiplo) para ECB y CBC, y la del texto plano para CFB y OFB. Se midió contra el oráculo para las 16 combinaciones (nota K5).

\textbf{Contraseña equivocada sin código de autenticación.} El formato no tiene MAC. Una contraseña, algoritmo o modo equivocados se detectan por el relleno PKCS5 (ECB y CBC) y por una validación estricta del encuadre del texto plano
(tamaño coherente, un NUL final, extensión imprimible), de modo que casi siempre hay un error claro (código de salida 2) y ningún archivo de salida. Es posible, aunque muy improbable, que una contraseña equivocada pase ambos controles y produzca
basura. Es una debilidad del formato, no del código, y se registró en la sección de modelo de amenazas de \texttt{docs/CRYPTO-NOTES.md}.

\textbf{Sin \texttt{-pass} no hay cifrado.} La consigna dice que sin contraseña se hace esteganografía sin cifrar. Si el usuario da \texttt{-a} o \texttt{-m} pero no \texttt{-pass}, el programa ignora esas opciones, y como el
silencio sería confuso imprime en el error estándar \texttt{stegobmp: note: -a/-m ignored because no -pass was given (no encryption)} (nota K8).

\subsection{Formato del mensaje y portador}

El tamaño va en cuatro bytes \emph{big-endian}, seguido de los datos, de la extensión (un punto y sus letras) y de un byte NUL; con cifrado se antepone el tamaño del texto cifrado. Como cualquier archivo estego puede haber sido
alterado, el tamaño leído se acota con la capacidad restante del portador antes de reservar memoria, y el texto plano descifrado se valida por completo antes de escribir nada (casos \texttt{decrypt-hostile-size-*} y
\texttt{cap-enc-sparse-5GiB}, resumen 03-02).

El portador también trae detalles. El bloque de píxeles se usa completo, con los bytes de relleno de fila incluidos y en el orden del archivo (no en el orden de la imagen); las filas van de abajo hacia arriba, aunque el programa admite
alturas negativas; el encabezado y cualquier byte posterior al bloque se copian sin cambios. Todas las salidas se escriben primero en un archivo temporal del mismo directorio y luego se renombran sobre el destino (\texttt{rename}), de modo que
un error nunca deja un archivo a medias, incluso si \texttt{-out} coincide con \texttt{-p} (documento \texttt{SKELETON.md} de la fase 1).

\subsection{La consigna y los archivos de ejemplo}

Los comandos de ejemplo de la consigna~\cite{catedra2026} no se pueden copiar y pegar. Están escritos con un guion largo (–) en lugar del guion común y con comillas tipográficas (“mensaje1.txt”): el texto extraído del PDF
contiene \texttt{–in}, \texttt{–p}, \texttt{–steg} y \texttt{–a}. Tecleado tal cual, \texttt{./stegobmp -embed –in x -p Ejemplo/lado.bmp -out y.bmp -steg LSB1} termina con código de salida 1 y el mensaje
\texttt{stegobmp: error: unknown parameter '–in'} seguido de las líneas de uso. Se optó por respetar la sintaxis exacta del enunciado y no aceptar el guion largo, porque la consigna exige respetar esa sintaxis. El ejemplo de extracción, además, escribe
\texttt{“imagenmas1 .bmp”} con un espacio antes de la extensión, que no es el nombre del archivo que se generó en el ejemplo anterior.

Otras incoherencias menores: el método se llama ``LSB Enhanced'' en la sección 5.1 y ``LSB Improved'' en el resto de la consigna, aunque el valor de \texttt{-steg} es siempre \texttt{LSBI}; el ejemplo 2 llama ``imagen'' a un archivo
\texttt{.txt}; los indicadores de patrones se guardan con los tres canales R, G y B mientras que el mensaje salta el rojo, cosa que sólo dice el paper (nota A2 y punto 16 de las incoherencias listadas en \texttt{docs/LSBI-NOTES.md}). Por último,
\texttt{Ejemplo/README.txt} no está en UTF-8: el comando \texttt{file} lo identifica como ``ISO-8859 text, with CRLF line terminators'', así que sus letras acentuadas aparecen como caracteres de reemplazo en una terminal UTF-8.

\subsection{Compilación y portabilidad}

La consigna exige que el programa compile y corra en Ubuntu 22.04.3 y en pampero, dos entornos con versiones distintas de OpenSSL (de la 1.0.2 a la 3.x en los sistemas que podían aparecer), cuya interfaz de \texttt{EVP} cambió de forma incompatible entre ellas
(por ejemplo, la creación y destrucción de contextos). El código se limita a llamadas comunes a 1.0.2, 1.1.1 y 3.x, y el caso \texttt{build-openssl-api-allowlist} de \texttt{make test} falla si aparece un identificador de OpenSSL fuera de una lista
permitida, un encabezado distinto de \texttt{evp.h}, \texttt{crypto.h} o \texttt{err.h}, o una biblioteca de enlace distinta de \texttt{libcrypto}.

Se verificó la compilación limpia y las pruebas en pampero (OpenSSL 3.6.3) y en Ubuntu 24.04 bajo WSL2 (OpenSSL 3.0.13, gcc 13.3). \textbf{El entorno exacto que nombra la consigna, Ubuntu 22.04.3, no se probó}: la distribución no estaba instalada en la máquina de desarrollo y se
renunció a esa verificación; lo que da confianza en su lugar es la lista de identificadores permitidos y las dos pruebas en entornos reales.

Tres problemas concretos de construcción quedaron documentados. Primero, \texttt{CC ?= gcc} en el \texttt{Makefile} no tiene efecto, porque \texttt{make} ya define \texttt{CC} como \texttt{cc}; se corrigió con
\texttt{ifeq (\$(origin CC),default)}, que conserva los cambios hechos desde la línea de comandos (resumen 03-03). Segundo, la dependencia de enlace se comprueba con las entradas \texttt{NEEDED} de \texttt{readelf} y no con
\texttt{ldd}, para que las dependencias transitivas de \texttt{libcrypto} en pampero no cuenten como dependencias del programa. Tercero, una copia de trabajo en Windows con \texttt{core.autocrlf=true} convierte los finales de línea en CRLF, con lo que
\texttt{bash} y \texttt{make} fallan bajo WSL; se resolvió con \texttt{.gitattributes}, que fuerza \texttt{eol=lf} en \texttt{*.c}, \texttt{*.h}, \texttt{*.sh}, \texttt{*.py} y \texttt{Makefile} y declara binarios los \texttt{.bmp}, \texttt{.png} y \texttt{.pdf}
para que los vectores nunca se conviertan como texto. Hay además fricciones propias de WSL: \texttt{/tmp} no se comparte entre invocaciones separadas de \texttt{wsl} (las comparaciones antes y después hubo que hacerlas dentro de una misma llamada) y, sobre el disco NTFS, \texttt{make} a veces
avisa \texttt{Clock skew detected}.

\subsection{Cómo se verificó}

\texttt{make test} ejecuta 242 casos (\texttt{SUMMARY: 242 passed, 0 failed}), y \texttt{make test-tools} ejecuta 42 casos (\texttt{SUMMARY: 42 passed, 0 failed}) de las herramientas de análisis. Incluyen: la extracción de los seis vectores de la cátedra y su reinserción byte a byte
(cada archivo se reproduce exactamente al insertar de nuevo el mismo mensaje con la misma clave); la matriz de métodos por algoritmos y modos, con un recorrido completo por LSB1, LSB4 y LSBI con y sin cifrado; casos de contraseña, algoritmo y modo equivocados, de
datos dañados y de tamaños hostiles; corridas con AddressSanitizer y UBSan; y la ejecución textual de los ejemplos del README por \texttt{tests/check\_readme.sh}, para que la documentación no pueda divergir del programa.

El descifrado se compara con un oráculo independiente, \texttt{tools/cryptoref.py}, que usa \texttt{hashlib.pbkdf2\_hmac} de Python y la línea de comandos de \texttt{openssl}, y no comparte ningún código con la implementación en C.
Esa independencia es la que permitió encontrar la errata del archivo de ejemplo. Los límites de la verificación están dichos arriba: AES-192 y varios modos (todas las combinaciones ECB, AES en CFB de 8 bits y 3DES en OFB) sólo están contrastados
con el oráculo y no con archivos de la cátedra, y los cuatro puntos de LSBI sin decidir se validarán cuando lleguen los archivos de estegoanálisis.
